%!TEX encoding = UTF-8 Unicode
\documentclass[11pt, a4paper]{article}

\usepackage[utf8]{inputenc}
\usepackage[T1]{fontenc}
\usepackage[french]{babel}
\usepackage{amsmath, amsthm, amssymb, amsfonts}
\usepackage[top=2cm, left=1.2cm, right=1.2cm]{geometry}
\usepackage{hyperref}
\usepackage{graphicx}
\usepackage{xcolor}
\usepackage{setspace}
\onehalfspacing
\usepackage{array}
\usepackage{booktabs}
\usepackage{float}
\usepackage{fancyhdr}
\usepackage{caption}
\usepackage{fixltx2e}
\usepackage{enumerate}

\usepackage{listings}
\usepackage{color}
\lstset{ %
  language=Matlab,                     % the language of the code
  basicstyle=\footnotesize,       % the size of the fonts that are used for the code
  numbers=left,                   % where to put the line-numbers
  numberstyle=\tiny\color{mygray},  % the style that is used for the line-numbers
  stepnumber=1,                   % the step between two line-numbers. If it's 1, each line
                                  % will be numbered
  numbersep=5pt,                  % how far the line-numbers are from the code
  backgroundcolor=\color{white},  % choose the background color. You must add \usepackage{color}
  showspaces=false,               % show spaces adding particular underscores
  showstringspaces=false,         % underline spaces within strings
  showtabs=false,                 % show tabs within strings adding particular underscores
  frame=single,                   % adds a frame around the code
  rulecolor=\color{black},        % if not set, the frame-color may be changed on line-breaks within not-black text (e.g. commens (green here))
  tabsize=2,                      % sets default tabsize to 2 spaces
  captionpos=b,                   % sets the caption-position to bottom
  breaklines=true,                % sets automatic line breaking
  breakatwhitespace=false,        % sets if automatic breaks should only happen at whitespace
  title=\lstname,                 % show the filename of files included with \lstinputlisting;
                                  % also try caption instead of title
  keywordstyle=\color{blue},      % keyword style
  commentstyle=\color{mygreen},   % comment style
  stringstyle=\color{mymauve},      % string literal style
  escapeinside={\%*}{*)},         % if you want to add a comment within your code
  morekeywords={*,...}            % if you want to add more keywords to the set
} 

\colorlet{punct}{red!60!black}
\definecolor{background}{HTML}{EEEEEE}
\definecolor{delim}{RGB}{20,105,176}
\colorlet{numb}{magenta!60!black}

\lstdefinelanguage{json}{
    basicstyle=\normalfont\ttfamily,
    numbers=left,
    numberstyle=\scriptsize,
    stepnumber=1,
    numbersep=8pt,
    showstringspaces=false,
    breaklines=true,
    frame=lines,
    backgroundcolor=\color{background},
    literate=
     *{0}{{{\color{numb}0}}}{1}
      {1}{{{\color{numb}1}}}{1}
      {2}{{{\color{numb}2}}}{1}
      {3}{{{\color{numb}3}}}{1}
      {4}{{{\color{numb}4}}}{1}
      {5}{{{\color{numb}5}}}{1}
      {6}{{{\color{numb}6}}}{1}
      {7}{{{\color{numb}7}}}{1}
      {8}{{{\color{numb}8}}}{1}
      {9}{{{\color{numb}9}}}{1}
      {:}{{{\color{punct}{:}}}}{1}
      {,}{{{\color{punct}{,}}}}{1}
      {\{}{{{\color{delim}{\{}}}}{1}
      {\}}{{{\color{delim}{\}}}}}{1}
      {[}{{{\color{delim}{[}}}}{1}
      {]}{{{\color{delim}{]}}}}{1},
}

\definecolor{mygreen}{rgb}{0,0.6,0}
\definecolor{mygray}{rgb}{0.5,0.5,0.5}
\definecolor{mymauve}{rgb}{0.58,0,0.82}


\renewcommand{\footrulewidth}{0.4pt}

%%%%%%% SETUP %%%%%%%%

% Header and footer
\setlength{\headheight}{15.2pt}
\pagestyle{fancy}
\lhead{Robotique }
\rhead{Documentation}
\chead{Asterix}

%%%% END OF SETUP %%%%

\begin{document}
%%%% Frontpage %%%%%%
\begin{titlepage}

\newcommand{\HRule}{\rule{\linewidth}{0.5mm}} % 
\center % Höfum allt á blaðsíðunni fyrir miðju

%----------------------------------------------------------------------------------------
% Forsiðumynd
%----------------------------------------------------------------------------------------

\begin{center}
\begin{minipage}{0.4\textwidth}
\end{minipage}\\[1cm]
%\vspace{2cm}
  % Upper part of the page
  \includegraphics[width=10cm]{logo.jpg}\\[1cm]
\end{center}
%----------------------------------------------------------------------------------------
% Heiti skóla
%----------------------------------------------------------------------------------------
 % Nafn skola
%\textsc{\Large Informatique}\\[0.25cm]
%Þetta:[xcm] segir til um bil fyrir neðan línu

%----------------------------------------------------------------------------------------
% Titill
%----------------------------------------------------------------------------------------

\HRule \\[0.4cm] %Lárétt lina
{ \huge \bfseries Introduction à la robotique par la pratique}\\[0.5cm] % Titill ritgerðar
\textsc{\Large Documentation de Projet}\\[0.5cm] 
\HRule \\[1.5cm]  %Lárétt lína
 
%------------------------------
% Höfundur
%------------------------------

\begin{minipage}{0.4\textwidth} %Setjum minipage með breidd 0.4 af breidd texta
\begin{flushleft} \large %Höfum þetta til vinstri
\emph{Étudiants:}
\\
Corentin Charles
\\
Clément Renazeau
\\
Þorsteinn Hjörtur Jónsson % Nafn Höfundar
\end{flushleft}
\end{minipage}
~
\begin{minipage}{0.4\textwidth}
\begin{flushright} \large %Höfum þetta til hægri
\vspace{-1cm}
\emph{Enseigneurs:} 
\\
Rémi Fabre
\\

 % Nafn Kennara
\end{flushright}
\end{minipage}\\[3cm]
%-----------------------------
%   Dagsetning
%-----------------------------
\nopagebreak      %Við skulum ekki fara á nýja blaðsíðu
{\large \today}\\[3cm] % Gott er að vita dagsetninguna
\null
\end{titlepage}
\clearpage
\section{Introduction}
This is the introduction.\\
\clearpage
\section{Rotation}
\subsection{Description}
\subsection{Program}
\lstinputlisting[language=python]{rotation.py}
\clearpage
\section{Walk}
\subsection{Description}
\subsection{Program}
\lstinputlisting[language=python]{walk.py}
\clearpage
\section{Main}
\subsection{Description}
\subsection{Program}
\lstinputlisting[language=python]{main.py}
\clearpage

\vspace{1.5cm}



\end{document}